\chapter{Sirkuit RC dan RL}\label{ch:g12:rc-rl-circuits}

Tekanlah tombol rana kamera setengah jalan: dengung tipis menanjak selama dua
sekon, sebuah lampu menyala hijau --- baru setelah itu lampu kilatnya menyala,
menumpahkan dalam semilisekon apa yang dikumpulkannya selama dua sekon. Bab ini
menemui dua komponen yang memberi sirkuit sebuah ingatan dan sebuah jam ---
\omterm{def:g12:rc-rl-circuits:capacitor}{kapasitor} dan \omterm{def:g11:magnetic-fields:solenoid}{kumparan} --- serta persamaan diferensial pertama dalam fisika,
yang dapat diperiksa dengan turunan tahun ini.

\section{Kapasitor}

\begin{definition}[Kapasitor dan kapasitansi]\label{def:g12:rc-rl-circuits:capacitor}
\emph{Kapasitor}\index{kapasitor} adalah sepasang keping logam yang berhadapan
dan dipisahkan oleh penyekat yang tipis. Ketika dikawatkan ke sumber, kepingnya
mengambil muatan yang berlawanan $+q$ dan $-q$ --- \omterm{def:g11:electric-gravitational-fields:efield}{medan listrik} memenuhi celahnya
(\cref{ch:g11:electric-gravitational-fields}) --- dan tegangan $u$
di antara kepingnya sebanding dengan muatan yang tersimpannya:
\[ q = C\,u . \]
Tetapan $C$ adalah \emph{kapasitansi}\index{kapasitansi}, dalam
\emph{farad}\index{farad} (\unit{F}): $\qty{1}{F} = \qty{1}{C/V}$.
\end{definition}

\begin{example}[Orde besaran]\label{ex:g12:rc-rl-circuits:orders}
Satu \omterm{def:g12:rc-rl-circuits:capacitor}{farad} itu raksasa --- satu \omterm{def:g11:fundamental-interactions:charge}{coulomb} terparkir tiap \omterm{def:g11:circuits-and-power:voltage}{volt}. Keramik pada sirkuit
radio: \unit{pF} sampai \unit{nF}; elektrolit pada catu \omterm{def:g11:work-of-force:power}{daya} dan
lampu kilat: \unit{\micro F} sampai \unit{mF}; superkapasitor: ribuan \omterm{def:g12:rc-rl-circuits:capacitor}{farad}.
\end{example}

\begin{proposition}[Arus yang masuk ke kapasitor]\label{prop:g12:rc-rl-circuits:cdudt}
Arus yang tiba pada keping \omterm{def:g12:rc-rl-circuits:capacitor}{kapasitor} adalah laju pertumbuhan
muatannya:
\[
i = \frac{dq}{dt} = C\,\frac{du}{dt} ,
\]
dengan huruf kecil menandai besaran yang berubah menurut waktu ($U$, $I$ tetap mantap,
\cref{ch:g11:circuits-and-power}).
\end{proposition}

\begin{proof}
Arus adalah muatan yang diserahkan tiap sekon
(\cref{def:g11:circuits-and-power:current}): pada selang yang menyusut,
$i$ adalah turunan $q$; dan $q = Cu$ dengan $C$ yang tetap memberikan
$C\,du/dt$.
\end{proof}

\begin{remark}[Dua akibatnya]\label{rem:g12:rc-rl-circuits:continuity}
Jika $u$ tetap, $i = 0$: \omterm{def:g12:rc-rl-circuits:capacitor}{kapasitor} yang bermuatan tidak melewatkan arus tetap.
Dan $u$ tidak pernah dapat meloncat --- loncatan akan menuntut arus tak hingga:
\omterm{def:g11:circuits-and-power:voltage}{tegangan} \omterm{def:g12:rc-rl-circuits:capacitor}{kapasitor} adalah ingatan sirkuitnya.
\end{remark}

\section{Mengisi kapasitor: sirkuit RC}

Satu simpal --- sumber $\mathcal{E}$ (\cref{def:g11:circuits-and-power:emf}),
sakelar, resistor, \omterm{def:g12:rc-rl-circuits:capacitor}{kapasitor} kosong: tutuplah $K$ pada $t = 0$.

\begin{omfigure}
\begin{circuitikz}[font=\small]
  \draw (0,0) to[battery1, l=$\mathcal{E}$] (0,2.4) to[nos, l=$K$] (2.1,2.4)
    to[R=$R$, i=$i$] (4.2,2.4) -- (5.4,2.4) to[C=$C$, v=$u$] (5.4,0) -- (0,0);
\end{circuitikz}

{\small Sirkuit pengisian RC: menutup $K$ membiarkan sumbernya mendorong muatan
melalui $R$ ke keping $C$.}
\end{omfigure}

\begin{theorem}[Hukum pengisian sirkuit RC]\label{thm:g12:rc-rl-circuits:charging}
\index{sirkuit RC}
Setelah sakelarnya tertutup, \omterm{def:g11:circuits-and-power:voltage}{tegangan} \omterm{def:g12:rc-rl-circuits:capacitor}{kapasitornya} menaati, lalu menyusul dari
$u(0) = 0$:
\[
RC\,\frac{du}{dt} + u = \mathcal{E},
\qquad
u(t) = \mathcal{E}\left(1 - e^{-t/\tau}\right),
\qquad \tau = RC .
\]
\end{theorem}

\begin{proof}
\omterm{def:g11:circuits-and-power:voltage}{Tegangan} sepanjang simpalnya berjumlah (\cref{prop:g11:circuits-and-power:laws}):
$\mathcal{E} = u_R + u = Ri + u$ pada setiap saat; sulihkanlah
$i = C\,du/dt$ (\cref{prop:g12:rc-rl-circuits:cdudt}). Kini \emph{periksalah}
calonnya: $du/dt = (\mathcal{E}/\tau)e^{-t/\tau}$, sehingga
$RC\,(\mathcal{E}/RC)\,e^{-t/\tau} + \mathcal{E}(1 - e^{-t/\tau}) =
\mathcal{E}$, yang benar untuk semua $t$; dan $u(0) = \mathcal{E}(1-1) = 0$. Bahwa
tidak ada kurva \emph{lain} yang cocok dengan persamaannya dan awalnya dibuktikan pada
pelajaran matematika tahun ini tentang persamaan diferensial.
\end{proof}

\begin{definition}[Tetapan waktu]\label{def:g12:rc-rl-circuits:tau}
\emph{Tetapan waktu}\index{tetapan waktu} sebuah \omterm{thm:g12:rc-rl-circuits:charging}{sirkuit RC} adalah
$\tau = RC$. \omterm{def:g10:orders-of-magnitude:unit}{Satuannya}: $\unit{\ohm} \times \unit{F} =
(\unit{V/A})(\unit{C/V}) = \unit{C/A} = \unit{s}$ --- irama milik
sirkuitnya sendiri.
\end{definition}

\begin{method}[Membaca $\tau$ pada kurvanya]\label{met:g12:rc-rl-circuits:reading}
Tiga bacaan yang setara, langsung dari grafiknya:
\begin{enumerate}
  \item \emph{garis singgung di titik asalnya:} $u'(0) = \mathcal{E}/\tau$, sehingga
        garis singgung awalnya mencapai asimtotnya pada $t = \tau$;
  \item \emph{63\%:} $u(\tau) = \mathcal{E}(1 - e^{-1}) \approx
        0.63\,\mathcal{E}$;
  \item \emph{95\%:} $u(3\tau) \approx 0.95\,\mathcal{E}$ --- terisi,
        untuk keperluan praktisnya ($5\tau$: $99\%$).
\end{enumerate}
\end{method}

\begin{proposition}[Pengosongan]\label{prop:g12:rc-rl-circuits:discharge}
\omterm{def:g12:rc-rl-circuits:capacitor}{Kapasitor} yang terisi sampai $U_0$ lalu ditutup pada resistor $R$ saja menaati
$RC\,du/dt + u = 0$ dan menyusul
\[
u(t) = U_0\,e^{-t/\tau}, \qquad \tau = RC
\]
--- $37\%$ tersisa pada $\tau$, $5\%$ pada $3\tau$.
\end{proposition}

\begin{proof}
Hukum simpal yang sama, tanpa sumbernya; penyulihan lagi:
$RC(-U_0/\tau)e^{-t/\tau} + U_0 e^{-t/\tau} = 0$, dengan $u(0) = U_0$.
\end{proof}

\begin{omfigure}
\begin{tikzpicture}
\begin{axis}[omaxis, width=9.2cm, height=5.4cm, xlabel={$t$}, ylabel={$u$},
    xmin=0, xmax=5, ymin=0, ymax=7.4, xtick={0,1,2,3,4,5},
    xticklabels={$0$,$\tau$,$2\tau$,$3\tau$,$4\tau$,$5\tau$},
    ytick={0,2.21,3.79,6.0},
    yticklabels={$0$,$0.37\,\mathcal{E}$,$0.63\,\mathcal{E}$,$\mathcal{E}$}]
  \addplot[omExo, dashed, domain=0:5, samples=2] {6};
  \addplot[omExo, dashed, domain=0:1.18, samples=2] {6*x};
  \addplot[omExo, dashed, domain=0:1.1, samples=2] {6-6*x};
  \addplot[omDef, thick, domain=0:5, samples=80] {6*(1-exp(-x))};
  \addplot[omThm, thick, domain=0:5, samples=80] {6*exp(-x)};
  \draw[dashed, omExo] (1,0) -- (1,3.79) -- (0,3.79);
  \draw[dashed, omExo] (3,0) -- (3,5.70);
\end{axis}
\end{tikzpicture}

{\small Pengisian (biru, dari $u = 0$ menuju $\mathcal{E}$) dan pengosongan
(merah, dari $U_0 = \mathcal{E}$). Kedua garis singgung awalnya menemui
asimtotnya pada $t = \tau$; $63\%$ rampung pada $\tau$, $95\%$ pada $3\tau$.}
\end{omfigure}

\begin{example}[Angka sebuah lampu kilat]\label{ex:g12:rc-rl-circuits:flashrc}
Sebuah lampu kilat foto mengisi $C = \qty{800}{\micro F}$ melalui
$R = \qty{1.0}{k\ohm}$: $\tau = 10^3 \times 8.0\times10^{-4} =
\qty{0.80}{s}$, sehingga lampu ``siap'' --- yang dikawatkan agar menyala pada $95\%$ ---
menunggu $3\tau \approx \qty{2.4}{s}$: dengung yang menggantung itu. Ketika ditembakkan
melalui tabung kilatnya yang hanya \qty{0.50}{\ohm}, \omterm{def:g12:rc-rl-circuits:capacitor}{kapasitor} yang \emph{sama}
mengosong dengan $\tau' = \qty{0.40}{ms}$: dua ribu kali
lebih cepat keluar daripada masuk.
\end{example}

\section{Kumparan: induktansi dan sirkuit RL}

\begin{definition}[Induktor dan induktansi]\label{def:g12:rc-rl-circuits:inductor}
\emph{Induktor}\index{induktor} adalah gulungan kawat: arusnya
menembus lilitannya sendiri dengan \omterm{def:g11:magnetic-fields:bfield}{medan magnet}
(\cref{ch:g11:magnetic-fields}); setiap kali arus itu \emph{berubah}, sebuah \omterm{def:g11:circuits-and-power:voltage}{tegangan}
muncul pada \omterm{def:g11:magnetic-fields:solenoid}{kumparannya}:
\[
u = L\,\frac{di}{dt} .
\]
Tetapan $L$ adalah \emph{induktansi}\index{induktansi}, dalam
\emph{henry}\index{henry} (\unit{H}): $\qty{1}{H} = \qty{1}{V.s/A}$.
\end{definition}
\admitted

\begin{remark}[Kumparan melawan perubahan]\label{rem:g12:rc-rl-circuits:oppose}
Mengapa arus yang berubah menghasilkan \omterm{def:g11:circuits-and-power:voltage}{tegangan} --- induksi --- diturunkan pada
jilid universitasnya; di sini kita menerima hukumnya beserta wataknya:
\omterm{def:g11:magnetic-fields:solenoid}{kumparannya} \emph{melawan perubahan} arusnya. Kembar cermin: $u$ sebuah \omterm{def:g12:rc-rl-circuits:capacitor}{kapasitor}
tidak dapat meloncat, $i$ sebuah \omterm{def:g11:magnetic-fields:solenoid}{kumparan} tidak dapat meloncat.
\end{remark}

\begin{omfigure}
\begin{circuitikz}[font=\small]
  \draw (0,0) to[battery1, l=$\mathcal{E}$] (0,2.4) to[nos, l=$K$] (2.1,2.4)
    to[R=$R$] (4.2,2.4) -- (5.4,2.4) to[L=$L$, i=$i$, v=$u$] (5.4,0) -- (0,0);
\end{circuitikz}

{\small \omterm{prop:g12:rc-rl-circuits:rl}{Sirkuit RL}: setelah $K$ tertutup, \omterm{def:g11:magnetic-fields:solenoid}{kumparannya} hanya membiarkan arusnya
menanjak pada iramanya sendiri $\tau = L/R$.}
\end{omfigure}

\begin{proposition}[Kenaikan arus pada sirkuit RL]\label{prop:g12:rc-rl-circuits:rl}
\index{sirkuit RL}
Menutup sakelarnya pada sumber $\mathcal{E}$, resistor $R$ dan \omterm{def:g11:magnetic-fields:solenoid}{kumparan}
$L$ yang terangkai seri memberikan $\mathcal{E} = Ri + L\,di/dt$, yaitu
\[
\frac{L}{R}\,\frac{di}{dt} + i = \frac{\mathcal{E}}{R},
\qquad\text{yang diselesaikan oleh}\quad
i(t) = \frac{\mathcal{E}}{R}\left(1 - e^{-t/\tau}\right),
\quad \tau = \frac{L}{R} .
\]
\end{proposition}

\begin{proof}
Hukum simpalnya lagi; dan inilah \emph{persamaan yang sama} seperti pada
\cref{thm:g12:rc-rl-circuits:charging} dengan $u \to i$,
$\mathcal{E} \to \mathcal{E}/R$, $RC \to L/R$ --- penyulihan yang sama
memeriksa penyelesaiannya, dan setiap bacaan pada
\cref{met:g12:rc-rl-circuits:reading} berpindah ke sini.
\end{proof}

\begin{example}[Sebuah elektromagnet terbangun]\label{ex:g12:rc-rl-circuits:rl}
\omterm{def:g11:magnetic-fields:solenoid}{Kumparan} sebuah relai, $L = \qty{0.30}{H}$ dan $R = \qty{6.0}{\ohm}$, pada catu
\qty{12}{V}: arus akhirnya $\mathcal{E}/R = \qty{2.0}{A}$, pada
irama $\tau = 0.30/6.0 = \qty{50}{ms}$ --- terjaga dalam $3\tau \approx
\qty{0.15}{s}$.
\end{example}

\begin{remark}[Bunga api saat dibuka]\label{rem:g12:rc-rl-circuits:spark}
Bukalah sakelar pada \omterm{def:g11:magnetic-fields:solenoid}{kumparan} yang sedang hidup dan $i$ harus runtuh ke nol dalam pecahan
milisekon: $di/dt$ menjadi raksasa dan $u = L\,di/dt$ mencapai ratusan
atau ribuan \omterm{def:g11:circuits-and-power:voltage}{volt} --- bunga api meloncati sentuhannya. \omterm{def:g11:magnetic-fields:solenoid}{Kumparan} pengapian sebuah mobil
melakukannya dengan sengaja, ribuan kali tiap menit.
\end{remark}

\section{Energi yang tersimpan}

\begin{proposition}[Energi di kapasitor, energi di kumparan]\label{prop:g12:rc-rl-circuits:energy}
\index{energi!kapasitor}\index{energi!kumparan}
\omterm{def:g12:rc-rl-circuits:capacitor}{Kapasitor} $C$ yang terisi sampai tegangan $u$ dan \omterm{def:g11:magnetic-fields:solenoid}{kumparan} $L$ yang membawa arus
$i$ menyimpan
\[
E_C = \tfrac12\,C u^2 = \frac{q^2}{2C},
\qquad
E_L = \tfrac12\,L i^2 .
\]
\end{proposition}

\begin{proof}
Mendaratkan muatan kecil $dq$ pada keping yang sudah bertegangan $u$ berongkos
$u\,dq$ (\cref{def:g11:circuits-and-power:voltage}); dengan menumpuk
ongkosnya, jumlahnya adalah luas di bawah garis $u = q/C$ --- segitiga
di bawah ini --- sehingga $E_C = \tfrac12 qu = \tfrac12 Cu^2$. Segitiga yang sama untuk
\omterm{def:g11:magnetic-fields:solenoid}{kumparannya} (menaikkan $i$ sebesar $di$ berongkos $Li\,di$) memberikan
$E_L = \tfrac12 Li^2$.
\end{proof}

\begin{omfigure}
\begin{tikzpicture}[font=\small, x=1.35cm, y=1.05cm]
  \fill[omDef!25] (0,0) -- (4,0) -- (4,2.6) -- cycle;
  \fill[omDef!60] (2.4,0) rectangle (2.7,1.658);
  \draw[omDef, thick] (0,0) -- (4,2.6);
  \draw[-Stealth] (0,0) -- (4.8,0) node[right] {$q$};
  \draw[-Stealth] (0,0) -- (0,3.1) node[above] {$u = q/C$};
  \draw[dashed] (4,0) -- (4,2.6) -- (0,2.6);
  \node[below] at (4,0) {$q_0$}; \node[left] at (0,2.6) {$u_0$};
  \node[below, font=\footnotesize] at (2.55,0) {$dq$};
  \node[omDef!50!black, font=\footnotesize] at (2.9,0.85)
    {$E_C$};
\end{tikzpicture}

{\small Tiap irisan $dq$ berongkos $u\,dq$; seluruh muatannya berongkos
luas segitiganya $E_C = \tfrac12 q_0 u_0$ --- \omterm{def:g11:fundamental-interactions:charge}{coulomb} yang lebih awal naik
pada \omterm{def:g11:circuits-and-power:voltage}{tegangan} yang rendah.}
\end{omfigure}

\begin{example}[Orde besaran energinya]\label{ex:g12:rc-rl-circuits:energies}
\omterm{def:g12:rc-rl-circuits:capacitor}{Kapasitor} lampu kilat pada \cref{ex:g12:rc-rl-circuits:flashrc}:
$\tfrac12 \times 8.0\times10^{-4} \times 300^2 = \qty{36}{J}$. Sebuah
superkapasitor \qty{3000}{F} pada \qty{2.7}{V}: \qty{1.1e4}{J} --- setara
sebuah baterai AA. \omterm{def:g11:magnetic-fields:solenoid}{Kumparan} \qty{10}{mH} pada \qty{10}{A}: \qty{0.50}{J}.
Kurang daripada yang disimpan baterai --- tetapi dapat dilepas dalam semilisekon.
\end{example}

\begin{remark}[Untuk apa kedua penyimpan waktu itu]\label{rem:g12:rc-rl-circuits:applications}
\index{penghalusan}\index{sirkuit pewaktu}
Masuk lambat, keluar cepat menyalakan lampu kilat kamera dan defibrilator.
Tundaan $\tau = RC$ menjalankan \omterm{rem:g12:rc-rl-circuits:applications}{sirkuit pewaktu}: penyapu kaca berselang,
suar berkedip, lampu tangga. \omterm{def:g12:rc-rl-circuits:capacitor}{Kapasitor} yang dipasang melintangi catu
\emph{menghaluskan} \omterm{def:g11:circuits-and-power:voltage}{tegangannya} ($u$ tidak dapat meloncat); \omterm{def:g11:magnetic-fields:solenoid}{kumparan} seri
menghaluskan arusnya ($i$ tidak dapat meloncat). Dan \omterm{def:g11:magnetic-fields:solenoid}{kumparan} beserta \omterm{def:g12:rc-rl-circuits:capacitor}{kapasitor}
bersama-sama membentuk \omterm{def:g12:oscillators-and-time:oscillator}{osilator} (\cref{ch:g12:rlc-oscillations}).
\end{remark}

\section{Latihan}

\begin{exercise}[$\star$]\label{exo:g12:rc-rl-circuits:1}
Sebuah \omterm{def:g12:rc-rl-circuits:capacitor}{kapasitor} \qty{470}{\micro F} diisi sampai \qty{12}{V}: berapa muatan yang
dipegangnya? \omterm{def:g11:circuits-and-power:voltage}{Tegangan} berapa yang akan meletakkan \qty{2.0}{mC} pada \qty{100}{\micro F}?
\end{exercise}

\begin{exercise}[$\star$]\label{exo:g12:rc-rl-circuits:2}
Bagikanlah \qty{47}{pF}, \qty{2200}{\micro F} dan \qty{10}{F} di antara sebuah
penghalus catu \omterm{def:g11:work-of-force:power}{daya}, sebuah penala radio dan sebuah superkapasitor; hitunglah \omterm{def:g10:energy-conservation:energy}{energi}
yang terakhir itu pada \qty{2.7}{V}.
\end{exercise}

\begin{exercise}[$\star$]\label{exo:g12:rc-rl-circuits:3}
\omterm{def:g11:circuits-and-power:voltage}{Tegangan} pada \omterm{def:g12:rc-rl-circuits:capacitor}{kapasitor} \qty{220}{\micro F} menanjak mantap dari $0$
sampai \qty{5.0}{V} dalam \qty{10}{ms}: berapa arus yang mengalir masuk? Berapa arusnya sekali
\omterm{def:g11:circuits-and-power:voltage}{tegangannya} diam, dan mengapa?
\end{exercise}

\begin{exercise}[$\star$]\label{exo:g12:rc-rl-circuits:4}
Hitunglah $\tau$ untuk $R = \qty{10}{k\ohm}$, $C = \qty{100}{\micro F}$, lalu
untuk $L = \qty{0.10}{H}$, $R = \qty{50}{\ohm}$. Periksalah, \omterm{def:g10:orders-of-magnitude:unit}{satuan} demi \omterm{def:g10:orders-of-magnitude:unit}{satuan}, bahwa
kedua gabungannya bersatuan sekon.
\end{exercise}

\begin{exercise}[$\star$]\label{exo:g12:rc-rl-circuits:5}
Sebuah \omterm{def:g12:rc-rl-circuits:capacitor}{kapasitor} terisi menuju \qty{6.0}{V}, melewati \qty{3.8}{V} pada
$t = \qty{2.0}{s}$. Bacalah \omterm{def:g12:rc-rl-circuits:tau}{tetapan waktunya}; kapan pengisiannya $95\%$ rampung?
\end{exercise}

\begin{exercise}[$\star\star$]\label{exo:g12:rc-rl-circuits:6}
Sebuah \omterm{thm:g12:rc-rl-circuits:charging}{sirkuit RC}: $\mathcal{E} = \qty{9.0}{V}$, $R = \qty{4.7}{k\ohm}$,
$C = \qty{220}{\micro F}$. Hitunglah $\tau$, $u(\tau)$, $u(3\tau)$, dan
arus pada saat sakelarnya tertutup.
\end{exercise}

\begin{exercise}[$\star\star$]\label{exo:g12:rc-rl-circuits:7}
Periksalah dengan penyulihan bahwa $u(t) = \mathcal{E}(1 - e^{-t/RC})$
memenuhi $RC\,du/dt + u = \mathcal{E}$; apa yang disandikan nilainya pada $t = 0$
dan limitnya untuk $t$ yang besar secara fisis?
\end{exercise}

\begin{exercise}[$\star\star$]\label{exo:g12:rc-rl-circuits:8}
Sebuah \omterm{def:g12:rc-rl-circuits:capacitor}{kapasitor} \qty{100}{\micro F} yang terisi sampai \qty{12}{V} mengosong melalui
\qty{47}{k\ohm}: hitunglah $\tau$, $u(\tau)$, dan waktu bagi $u$ untuk turun
di bawah \qty{0.60}{V}.
\end{exercise}

\begin{exercise}[$\star\star$]\label{exo:g12:rc-rl-circuits:9}
Arus pada \omterm{def:g11:magnetic-fields:solenoid}{kumparan} \qty{0.50}{H} menanjak mantap dari $0$ sampai
\qty{2.0}{A} dalam \qty{40}{ms}: \omterm{def:g11:circuits-and-power:voltage}{tegangan} berapa yang muncul padanya? Dengan
$u = L\,di/dt$, mengapa \emph{membuka} sakelar pada \omterm{def:g11:magnetic-fields:solenoid}{kumparan} menimbulkan bunga api, bukan menutupnya?
\end{exercise}

\begin{exercise}[$\star\star$]\label{exo:g12:rc-rl-circuits:10}
Relai pada \cref{ex:g12:rc-rl-circuits:rl}: periksalah dengan penyulihan bahwa
$i(t) = (\mathcal{E}/R)(1 - e^{-Rt/L})$ memenuhi
$L\,di/dt + Ri = \mathcal{E}$, lalu hitunglah $i(\tau)$ dan $i(3\tau)$.
\end{exercise}

\begin{exercise}[$\star\star$]\label{exo:g12:rc-rl-circuits:11}
Tunjukkan bahwa $u'(0) = \mathcal{E}/\tau$, lalu perolehlah bahwa garis singgung di
titik asalnya menemui asimtotnya pada $t = \tau$. Pada kurva terekam yang
memberikan $\tau = \qty{0.50}{s}$, dengan $R = \qty{2.0}{k\ohm}$: carilah $C$.
\end{exercise}

\begin{exercise}[$\star\star\star$]\label{exo:g12:rc-rl-circuits:12}
Urutkanlah menurut \omterm{def:g10:energy-conservation:energy}{energi} yang tersimpannya, dengan menghitung masing-masingnya: superkapasitor
\qty{3000}{F} pada \qty{2.7}{V}; \omterm{def:g12:rc-rl-circuits:capacitor}{kapasitor} lampu kilat, \qty{800}{\micro F} pada
\qty{300}{V}; \omterm{def:g11:magnetic-fields:solenoid}{kumparan} \qty{10}{mH} pada \qty{10}{A}; sebuah baterai AA
($E \approx \qty{1.3e4}{J}$). Apa yang ditawarkan ketiga yang pertama dan tidak dapat
ditawarkan baterainya?
\end{exercise}

\begin{exercise}[$\star\star\star$]\label{exo:g12:rc-rl-circuits:13}
Sebuah pewaktu lampu tangga menyala selama $u < 0.63\,\mathcal{E}$ pada
\omterm{def:g12:rc-rl-circuits:capacitor}{kapasitor} pengisiannya, lalu padam pada $t = \tau$. Rancanglah lampunya agar bertahan
\qty{5.0}{s} dengan $C = \qty{100}{\micro F}$: berapa $R$? Untuk ambang
$0.50\,\mathcal{E}$, tunjukkan bahwa tundaannya $\tau \ln 2$ lalu hitunglah.
\end{exercise}

\begin{exercise}[$\star\star\star$]\label{exo:g12:rc-rl-circuits:14}
Sebuah \omterm{def:g12:rc-rl-circuits:capacitor}{kapasitor} \qty{2200}{\micro F} menghaluskan sebuah catu: di antara pengisian ulangnya,
ia sendirian menyuapi beban yang menarik \qty{0.50}{A} selama \qty{10}{ms}. Dari
$i = C\,du/dt$, hitunglah lendutan \omterm{def:g11:circuits-and-power:voltage}{tegangannya}, lalu \omterm{def:g12:rc-rl-circuits:capacitor}{kapasitansi} yang
menjaganya di bawah \qty{0.50}{V}.
\end{exercise}

\begin{exercise}[$\star\star\star$]\label{exo:g12:rc-rl-circuits:15}
Sebuah \omterm{def:g11:magnetic-fields:solenoid}{kumparan} pengapian, $L = \qty{0.80}{H}$, membawa \qty{2.0}{A} ketika
pemutusnya memotongnya dalam sekitar \qty{1.0}{ms}. Taksirlah \omterm{def:g11:circuits-and-power:voltage}{tegangan} \omterm{def:g11:magnetic-fields:solenoid}{kumparannya}
selama pemotongan itu dan \omterm{def:g10:energy-conservation:energy}{energi} bunga apinya; mengapa pemutusan yang lambat
tidak memberikan bunga api?
\end{exercise}

\section{Soal: Lampu kilat kamera}

\begin{problem}[{Soal akhir pekan --- lampu kilat kamera: dengung dua
sekon, semilisekon kemuliaan, dan sepupunya yang terpasang di dinding yang
menghidupkan kembali jantung --- satu eksponensial menjalankan semuanya}]
\label{pb:g12:rc-rl-circuits:1}
Sebuah lampu kilat ringkas menyimpan \omterm{def:g10:energy-conservation:energy}{energi} di dalam \omterm{def:g12:rc-rl-circuits:capacitor}{kapasitor}
$C = \qty{800}{\micro F}$, yang diisi melalui $R = \qty{1.0}{k\ohm}$ dari
sebuah \omterm{def:g10:energy-conservation:transfer}{pengubah} yang dimodelkan sebagai sumber ideal $\mathcal{E} = \qty{300}{V}$;
lampu ``siap'' menyala pada $95\%$ \omterm{def:g11:circuits-and-power:voltage}{tegangan} penuhnya. Ketika ditembakkan,
\omterm{def:g12:rc-rl-circuits:capacitor}{kapasitornya} hanya melihat tabung kilatnya, $R_f = \qty{0.50}{\ohm}$. Datanya: sebuah
sel AA menyerahkan \qty{1.5}{V} dan \qty{2.4}{A.h}; $e = \qty{1.60e-19}{C}$.

\medskip
\noindent\textbf{Bagian I --- Wadah simpanannya.}
\begin{enumerate}
  \item Hitunglah muatan pada \omterm{def:g12:rc-rl-circuits:capacitor}{kapasitornya} pada \qty{300}{V}.
  \item Berapa \omterm{def:g11:fundamental-interactions:elementary}{muatan dasar} banyaknya itu?
  \item Hitunglah \omterm{def:g10:energy-conservation:energy}{energi} yang tersimpannya $E_C$.
  \item Sel AA-nya memegang $q = It \approx \qty{8.6e3}{C}$, jadi sekitar
        \qty{1.3e4}{J}. Berapa kali \omterm{def:g10:energy-conservation:energy}{energi} \omterm{def:g12:rc-rl-circuits:capacitor}{kapasitornya} nilai itu?
  \item Lalu mengapa harus bersusah payah dengan \omterm{def:g12:rc-rl-circuits:capacitor}{kapasitornya}? Satu kalimat, dengan kata ``\omterm{def:g11:work-of-force:power}{daya}''.
\end{enumerate}

\medskip
\noindent\textbf{Bagian II --- Menanti hijaunya.}
\begin{enumerate}[resume]
  \item Tulislah hukum simpalnya lalu ubahlah menjadi persamaan diferensial
        untuk $u(t)$.
  \item Periksalah dengan penyulihan bahwa $u(t) = \mathcal{E}(1 - e^{-t/RC})$
        menyelesaikannya, bermula dari $u(0) = 0$.
  \item Hitunglah \omterm{def:g12:rc-rl-circuits:tau}{tetapan waktunya} $\tau$.
  \item Setelah waktu berapa lampu siapnya menyala hijau, dan pada \omterm{def:g11:circuits-and-power:voltage}{tegangan}
        berapa? Kenalilah dengung dua sekonnya.
  \item Hitunglah arus pengisian awalnya; mengapa arus itu lalu meredup
        hilang?
\end{enumerate}

\medskip
\noindent\textbf{Bagian III --- Semilisekon kemuliaan.}
\begin{enumerate}[resume]
  \item Tulislah hukum pengosongannya $u(t)$ melalui tabungnya lalu hitunglah
        \omterm{def:g12:rc-rl-circuits:tau}{tetapan waktu} barunya $\tau'$.
  \item Hitunglah arus puncak melalui tabungnya.
  \item Hitunglah \omterm{def:g11:work-of-force:power}{daya} puncaknya; bandingkan dengan ceret \qty{2.2}{kW}.
  \item Dengan menganggap kilatnya usai pada $3\tau'$, sebutkan lamanya dan
        \omterm{def:g11:work-of-force:power}{daya} rata-ratanya.
  \item Bentuklah nisbah waktu pengisiannya terhadap waktu kilatnya: \omterm{def:g12:rc-rl-circuits:capacitor}{kapasitor} yang sama,
        muatan yang sama --- apa yang berubah, dan itu menjadikan \omterm{def:g12:rc-rl-circuits:capacitor}{kapasitor}
        sebuah tangki \omterm{def:g10:energy-conservation:energy}{energi}, atau sebuah tuas \omterm{def:g11:work-of-force:power}{daya}?
\end{enumerate}

\medskip
\noindent\textbf{Bagian IV --- Sepupunya di dinding.}
Sebuah defibrilator harus menyerahkan $E_C = \qty{200}{J}$ pada
$U_0 = \qty{1500}{V}$; dada dan pelatnya menyajikan sekitar \qty{50}{\ohm}.
\begin{enumerate}[resume]
  \item Hitunglah \omterm{def:g12:rc-rl-circuits:capacitor}{kapasitansi} yang dibutuhkannya.
  \item Hitunglah muatan yang tersimpannya.
  \item Hitunglah \omterm{def:g12:rc-rl-circuits:tau}{tetapan waktu} dan lama kejutannya ($3\tau$);
        bandingkan dengan lampu kilatnya.
  \item Hitunglah arus puncak dan \omterm{def:g11:work-of-force:power}{daya} puncak pada pasiennya.
  \item Sirkuit pengisian ulangnya harus membuat alatnya siap ($95\%$) dalam
        \qty{6.0}{s}: carilah \omterm{def:g11:circuits-and-power:resistance}{hambatan} pengisian yang dibutuhkannya dan
        arus pengisian awalnya --- angka yang akan dipesan perancangnya.
\end{enumerate}
\end{problem}
